\subsection*{A.1 Reproduction}

All results are deterministic given the seeds. The pipeline has three scripts.

\begin{enumerate}[leftmargin=1.8em,itemsep=2pt]
  \item \texttt{q1\_q2.py} --- descriptive statistics, Kruskal--Wallis and Mann--Whitney tests
        with Holm correction, repeated stratified five-fold cross-validation, date-holdout
        validation, composite fidelity, monotone fitting, bootstrap intervals and the Weibull
        comparisons of Problem~2. Seed 2026.
  \item \texttt{generate.py} --- paired McNemar tests, GEE models, the eight-combination
        permutation test, the two-part AFT model, the nine-row prediction table, CEI curves and
        radii, expected responders and the scenario analysis. Block bootstrap: 2000 resamples of
        the 30 blocks; combination permutation: 5000 swaps.
  \item \texttt{make\_figs.py} --- the eight figures.
\end{enumerate}

Environment: Python 3.14 with numpy 2.5, pandas 3.0, matplotlib 3.11 and statsmodels for the GEE
and AFT fits; XeLaTeX for typesetting. Total runtime under five minutes.

\subsection*{A.2 Data and code availability}

The data are those distributed with the 2026 APMCM Problem C and archived by the problem setters
in the public Zenodo record cited in the problem statement (Data01--11 and Supplementary~1--5).
No external data were used. The three scripts above read only the nine correlation files, Data01
and Data02 from the supplied archive; they, together with the intermediate result tables, are
provided as supplementary material, so every number in this paper can be regenerated exactly.


\subsection*{A.3 Date-adjusted sensitivity for the call-type effect}

Date is strongly associated with call type (Cram\'er's $V=0.528$), so the call-type odds ratios
could in principle be an artefact of when each call type was played. Table~\ref{tab:date-adj}
refits the model of Equation~\eqref{eq:gee} with a fixed effect for each recording date
(12 levels, 19 parameters, 139 clusters).

\begin{table}[htbp]
  \centering
  \caption{Call-type odds ratios before and after adjusting for recording date.}\label{tab:date-adj}
  \small
  \begin{tabular}{lrrrr}
    \toprule
    & \multicolumn{2}{c}{Unadjusted} & \multicolumn{2}{c}{Date-adjusted} \\
    \cmidrule(lr){2-3}\cmidrule(lr){4-5}
    Call type & OR [95\% CI] & $p$ & OR [95\% CI] & $p$ \\
    \midrule
    flat & 13.75 [6.58, 28.73] & $<0.001$ & 15.79 [6.44, 38.73] & $<0.001$ \\
    mod & 16.28 [6.39, 41.47] & $<0.001$ & 45.98 [13.13, 161.07] & $<0.001$ \\
    tv & 22.38 [9.47, 52.89] & $<0.001$ & 38.73 [12.90, 116.26] & $<0.001$ \\
    \bottomrule
  \end{tabular}
\end{table}

Adjustment increases the modulated and two-voice odds ratios by a factor of two to three and
widens their intervals substantially, but every effect remains significant. The date-confound
therefore inflates our uncertainty rather than explaining the effect away.

\subsection*{A.4 The recruitment model behind the index}

The index of Equation~\eqref{eq:cei} uses a second response definition, recruitment
(approach or vocalisation), fitted with the same estimating equations. Table~\ref{tab:gee-recruit}
reports it. Its most important entry is the distance term.

\begin{table}[htbp]
  \centering
  \caption{GEE for recruitment (approach or vocalisation). n = 278, 139 clusters.}\label{tab:gee-recruit}
  \small
  \begin{tabular}{lrrr}
    \toprule
    Term & OR & 95\% CI & $p$ \\
    \midrule
    flat vs control & 54.66 & 16.32--183.07 & $<0.001$ \\
    mod vs control & 42.59 & 13.08--138.63 & $<0.001$ \\
    tv vs control & 68.04 & 20.46--226.23 & $<0.001$ \\
    $\log$ distance (per $e$-fold) & 0.873 & 0.529--1.439 & 0.593 \\
    Group size (per bird) & 0.944 & 0.898--0.992 & 0.023 \\
    Context: commuting vs after & 0.794 & 0.340--1.857 & 0.595 \\
    Context: searching vs after & 2.677 & 0.891--8.044 & 0.079 \\
    \bottomrule
  \end{tabular}
\end{table}

Recruitment declines only weakly with distance and not significantly ($p=0.59$), whereas the
any-response outcome of Table~\ref{tab:q3-gee} declines clearly ($p=0.008$). The difference is
informative: distance mainly suppresses the cheap responses (heads-up, vocalisation) rather than
the decision to approach. It also means that the distance profile of the index is driven by
fidelity rather than by behaviour.

\subsection*{A.5 Software and source code}

Software: Python 3.14 with numpy 2.5, pandas 3.0, matplotlib 3.11 and statsmodels 0.15;
XeLaTeX for typesetting.

The complete source code is provided in the supplementary archive submitted with this paper
(\texttt{Capmcm26XXXXXXfj.zip}), as required by the submission specification. It consists of
three files, listed here with their size and the results they produce.

\begin{table}[htbp]
  \centering
  \caption{Source files in the supplementary archive.}\label{tab:code}
  \small
  \begin{tabular}{llrl}
    \toprule
    File & Produces & Lines & Random seed \\
    \midrule
    \texttt{q1\_q2.py} & Problem 1 and Problem 2 results, Figures 1--2 & 473 & 2026 \\
    \texttt{generate.py} & Problem 3 and Problem 4 results, Tables 8--22 & 391 & fixed \\
    \texttt{make\_figs.py} & Figures 1--8 & 232 & fixed \\
    \bottomrule
  \end{tabular}
\end{table}

The scripts read only the data files distributed with the problem and write plain CSV tables, so
every number in this paper can be regenerated without any private dependency. The code is not
reproduced in the appendix because its string literals contain the Chinese names of the output
tables, which would violate the English-only requirement for the paper.
